\begin{table}[htbp]
\centering
\small
\caption{Distribuzione e importi per città della stazione appaltante}
\begin{tabular}{lrrrr}
\toprule
Città & CIG & Percentuale & Importo complessivo & Mediana \\
\midrule
Roma & 3 & 14,3\% & 950.921,50 € & 400.000,00 € \\
Torino & 2 & 9,5\% & 752.254,34 € & 376.127,17 € \\
Carlino & 1 & 4,8\% & 45.001,02 € & 45.001,02 € \\
Grosseto & 1 & 4,8\% & 76.216,47 € & 76.216,47 € \\
Cremona & 1 & 4,8\% & 4.918,03 € & 4.918,03 € \\
Firenze & 1 & 4,8\% & 319.800,00 € & 319.800,00 € \\
Latina & 1 & 4,8\% & 137.493,90 € & 137.493,90 € \\
Spoleto & 1 & 4,8\% & 96.150,00 € & 96.150,00 € \\
Milano & 1 & 4,8\% & 75.000,00 € & 75.000,00 € \\
San Fratello & 1 & 4,8\% & 43.800,00 € & 43.800,00 € \\
Frascati & 1 & 4,8\% & 872.727,27 € & 872.727,27 € \\
Pomezia & 1 & 4,8\% & 48.000,00 € & 48.000,00 € \\
Castiglione Cosentino & 1 & 4,8\% & 715.471,24 € & 715.471,24 € \\
Colleferro & 1 & 4,8\% & 76.500,00 € & 76.500,00 € \\
Torriana & 1 & 4,8\% & 60.000,00 € & 60.000,00 € \\
Napoli & 1 & 4,8\% & 80.457,59 € & 80.457,59 € \\
Bressanone & 1 & 4,8\% & 84.236,15 € & 84.236,15 € \\
Verona & 1 & 4,8\% & 66.000,00 € & 66.000,00 € \\
\bottomrule
\end{tabular}
\end{table}
